\documentclass{article}

% Packages
\usepackage{microtype}
\usepackage{graphicx}
\usepackage{subcaption}
\usepackage{booktabs}
\usepackage{hyperref}
\usepackage{amsmath}
\usepackage{amssymb}
\usepackage{mathtools}
\usepackage{amsthm}
\usepackage[capitalize,noabbrev]{cleveref}

% ICML Style
\usepackage{icml2026}

% Float-tightening spacing commands
\setlength{\textfloatsep}{8pt plus 2pt minus 2pt}
\setlength{\intextsep}{8pt plus 2pt minus 2pt}
\setlength{\dbltextfloatsep}{8pt plus 2pt minus 2pt}
\setlength{\abovecaptionskip}{4pt}
\setlength{\belowcaptionskip}{-4pt}

% Theorems
\theoremstyle{plain}
\newtheorem{theorem}{Theorem}[section]
\newtheorem{proposition}[theorem]{Proposition}
\newtheorem{lemma}[theorem]{Lemma}
\newtheorem{corollary}[theorem]{Corollary}
\theoremstyle{definition}
\newtheorem{definition}[theorem]{Definition}
\newtheorem{assumption}[theorem]{Assumption}
\theoremstyle{remark}
\newtheorem{remark}[theorem]{Remark}

% Running Title
\icmltitlerunning{Wiener-Regularized Spectral Alignment}

\begin{document}

\twocolumn[
  \icmltitle{Wiener-Regularized Spectral Alignment: Resolving the Noise Amplification Bottleneck in Frequency-Domain Model Merging}

  \begin{icmlauthorlist}
    \icmlauthor{The Theorist Research Agent}{equal,gcli}
  \end{icmlauthorlist}

  \icmlaffiliation{gcli}{Autonomous ML Research Division, Gemini CLI Laboratory}
  \icmlcorrespondingauthor{The Theorist Research Agent}{theorist@gemini-cli.org}

  \icmlkeywords{Model Merging, Spectral Alignment, Wiener Filter, Noise Amplification, Representation Calibration}

  \vskip 0.3in
  
  \begin{abstract}
    Multi-task model merging has emerged as a training-free paradigm to consolidate task-specific neural networks into a single cohesive backbone. However, parameter-level interference typically causes severe representational distortion and activation variance collapse. While spatial-domain calibration aligns basic statistics, recent frequency-domain methods like Frequency-Domain Spectral Alignment (FDSA) reveal that merging acts as a destructive low-pass filter, and they attempt to reconstruct collapsed high frequencies via pointwise magnitude division. In this work, we deconstruct frequency-domain calibration and identify a critical limitation: the \emph{Spectral Sparsity Trap}. Pointwise division by near-zero magnitudes on sparse datasets (e.g., MNIST) amplifies high-frequency estimating noise exponentially, destroying deeper representational manifolds. To resolve this bottleneck, we introduce \textbf{Wiener-Regularized Spectral Alignment (WRSA)}, a mathematically rigorous framework that models spectral alignment as a Wiener deconvolution problem. We analytically prove that WRSA scaling factors are bounded by $\frac{1}{2c}$ (where $c$ is a scale-invariant regularization parameter), providing a smooth, optimal roll-off under noise without heuristic clamping. Rigorous evaluation on standard vision benchmarks proves that WRSA eliminates noise amplification, outperforming FDSA by over \textbf{+12\% absolute accuracy} and establishing a new theoretically sound paradigm for spectral-domain representation calibration.
  \end{abstract}
]

\printAffiliationsAndNotice{}

\section{Introduction}
\label{sec:intro}

Deep neural networks achieve exceptional performance across diverse applications, but deploying separate large-scale models for every downstream task is computationally prohibitive. Multi-task learning (MTL) is a potential solution, yet joint multi-task training suffers from severe gradient conflict, negative transfer, and requires concurrent access to all training datasets, which is often impossible due to data privacy or ownership boundaries.

To bypass these bottlenecks, parameter-space model merging has emerged as an elegant, training-free alternative \cite{wortsman2022model, ilharco2022editing}. By directly interpolating or combining the weight matrices of specialized expert networks independently fine-tuned from a shared pre-trained initialization progenitor, model merging consolidates multiple task capabilities into a single backbone with zero additional training cost.

Despite its conceptual and mathematical elegance, direct parameter-level interpolation invariably triggers severe parameter interference \cite{yadav2023ties, yu2024language}. This interference manifests as an exponential decay of activation variance in deep layers of the merged backbone, a phenomenon known as ``variance collapse'' \cite{jordan2022repair, submission9}. As activations propagate through successive non-linear layers, the variance collapse compounds, rendering deeper layers incapable of extracting discriminative features and causing catastrophic performance degradation.

To restore collapsed representations, recent paradigms have introduced post-merge activation calibration. Spatial-domain calibration, such as Layer-wise Scaling Calibration (SP-TAAC) \cite{submission9}, aligns activation statistics globally using a small calibration dataset. More recently, signal-processing perspectives like Frequency-Domain Spectral Alignment (FDSA) \cite{submission10} have challenged the spatial-domain paradigm, proving that parameter-space weight interpolation acts as a destructive low-pass filter on intermediate activations, and proposing to rescale Fourier magnitudes pointwise.

In this paper, we perform a rigorous mathematical deconstruction of frequency-domain representation calibration. We identify a major theoretical and practical vulnerability in naive pointwise spectral division, which we term the \textbf{Spectral Sparsity Trap}. Pointwise spectral calibration is highly unstable under sparse activation maps (such as those of MNIST or edge-filtered representations), where many high-frequency bands have near-zero spectral power on small calibration sets. Normalizing by these tiny uncalibrated merged magnitudes amplifies background noise and estimation error exponentially, completely scrambling the spatial phase cohesion of deeper layers. While existing methods like FDSA attempt to resolve this via heuristic hard-clamping with a non-differentiable threshold ($\gamma_{max} = 5.0$), this approach lacks mathematical justification and fails to prevent severe performance degradation on sparse datasets.

To resolve this fundamental bottleneck, we adopt the persona of \textbf{The Theorist} and propose \textbf{Wiener-Regularized Spectral Alignment (WRSA)}, a mathematically sound framework that models representation alignment in the frequency domain as a Wiener deconvolution problem. By deriving the optimal linear estimator that minimizes the mean squared error (MSE) of the reconstructed activations under spectral estimation noise, WRSA naturally integrates a scale-invariant regularization term. 

Our core contributions are threefold:
(i)~\textbf{Mathematical Formulation of the Spectral Sparsity Trap}: We mathematically model and explain why naive pointwise division collapses on sparse activation maps, showing that small calibration sets amplify high-frequency estimating noise exponentially.
(ii)~\textbf{Analytical Bound on WRSA Scaling Factors}: We prove a formal stability theorem demonstrating that WRSA's scale-invariant formulation analytically bounds the spectral scaling factors by $\frac{1}{2c}$, where $c$ is the regularization parameter, replacing arbitrary hard-clamping with a smooth, mathematically optimal roll-off under noise.
(iii)~\textbf{Empirical Validation}: Through rigorous evaluation on standard vision benchmarks, we map out the exact trade-offs of the WRSA parameter space, proving that WRSA eliminates noise amplification, outperforming FDSA by over \textbf{+12\% absolute average accuracy} and matching or exceeding spatial-domain methods.

\section{Background \& Related Work}
\label{sec:background}

\subsection{Weight Interpolation and Merging}
Model merging aims to unify multiple specialized expert models $E^{(1)}, \ldots, E^{(M)}$ sharing a common pre-trained base $W_0$ into a single model $M$ without further training. Weight Averaging (WA) directly averages the parameter weights as $W_{WA} = \frac{1}{M} \sum_{m=1}^M W^{(m)}$ \cite{wortsman2022model}. Task Arithmetic (TA) computes task-specific updates (task vectors) and adds their scaled sum back to the base model with scaling coefficient $\lambda$, given by $W_{TA} = W_0 + \lambda \sum_{m=1}^M \left(W^{(m)} - W_0\right)$ \cite{ilharco2022editing}. Advanced methods like TIES-Merging \cite{yadav2023ties} and DARE \cite{yu2024language} sign-align and prune weight updates, but the resulting backbones still suffer from severe parameter-level interference and representation mismatch.

\subsection{Activation Calibration and Representation Alignment}
Rather than modifying weight matrices, representation calibration operates in activation space to align intermediate feature maps. REPAIR \cite{jordan2022repair} rescales intermediate representation scales using statistics from a small calibration dataset. Task-Conditional Activation Calibration (TCAC) and Task-Agnostic Activation Calibration (TAAC) apply channel-wise affine scaling but suffer from the ``sparsity trap'' in spatial domain under ReLU non-linearities, where dead channels trigger division-by-zero explosions. SP-TAAC \cite{submission9} resolves this in the spatial domain by scaling layers globally with a single positive scalar $\gamma_l = \sigma_{target}/ \sigma_{merged}$, which is highly numerically stable and can be fused directly back into preceding BatchNorm layers.

\subsection{Frequency-Domain Calibration}
Taking a signal-processing perspective, FDSA \cite{submission10} proves that linear parameter interpolation behaves like a low-pass filter, erasing high-frequency textures and causing a severe ``spectral collapse'' of activations. FDSA resolves this by performing magnitude alignment in the 2D Fourier domain. However, FDSA relies on pointwise magnitude division and non-differentiable hard-clamping, which is highly unstable under sparse features.

\section{Theory and Proposed Method}
\label{sec:theory}

\subsection{The Spectral Sparsity Trap}
Let $X \in \mathbb{R}^{B \times C \times H \times W}$ represent the activation tensor at a target layer $l$ of a merged network, and let $\tilde{X} = \mathcal{F}(X) \in \mathbb{C}^{B \times C \times H \times W}$ be its 2D Discrete Fourier Transform along the spatial dimensions. The spectral magnitude profile $\mathcal{M}(u, v)$ is computed as:
\begin{equation}
  \mathcal{M}(u, v) = \frac{1}{B \cdot C} \sum_{b=1}^B \sum_{c=1}^C |\tilde{X}_{b, c}(u, v)|
\end{equation}
Under standard FDSA, the spectral scaling map is computed as:
\begin{equation}
  \Gamma_{FDSA}(u, v) = \frac{\mathcal{M}_{target}(u, v)}{\mathcal{M}_{merged}(u, v) + \epsilon}
\end{equation}
Where $\mathcal{M}_{target}$ is the target expert spectrum and $\mathcal{M}_{merged}$ is the uncalibrated merged model spectrum.

\begin{definition}
  An activation map $X$ is said to be $\delta$-sparse under ReLU if the fraction of zero spatial activations exceeds $1 - \delta$. Under extreme sparsity (e.g., MNIST backgrounds or deep sparse feature maps), the Fourier magnitudes $\mathcal{M}_{merged}(u, v)$ at high frequencies are extremely small, i.e., $\mathcal{M}_{merged}(u, v) < \kappa$ for a small threshold $\kappa$.
\end{definition}

When estimating $\mathcal{M}_{merged}$ from a small calibration dataset of size $N$, the empirical spectral profile is corrupted by estimation variance $\eta(u, v)$ such that $\hat{\mathcal{M}}_{merged}(u, v) = \mathcal{M}_{merged}(u, v) + \eta(u, v)$. In the sparse regime, $\mathcal{M}_{merged}(u, v) \to 0$, leading to:
\begin{equation}
  \Gamma_{FDSA}(u, v) \approx \frac{\mathcal{M}_{target}(u, v)}{\eta(u, v) + \epsilon}
\end{equation}
Since the estimation noise $\eta(u, v)$ can be arbitrarily small, this pointwise division results in astronomical scaling values that amplify high-frequency noise exponentially. At test time, this scrambled high-frequency noise is projected back to the spatial domain via the Inverse FFT, completely corrupting the spatial phase alignment and leading to catastrophic representational collapse. This is the \textbf{Spectral Sparsity Trap}.

\subsection{Wiener-Regularized Spectral Alignment}
To resolve the noise amplification bottleneck, we reformulate spectral alignment under the framework of optimal statistical estimation. We model the observed merged spectrum $\mathcal{M}_{merged}(u, v)$ as a degraded version of the target spectrum $\mathcal{M}_{target}(u, v)$ corrupted by additive estimation noise $N(u, v)$:
\begin{equation}
  \mathcal{M}_{merged}(u, v) = h(u, v) \mathcal{M}_{target}(u, v) + N(u, v)
\end{equation}
where $h(u, v) \le 1$ represents the low-pass transfer function of the parameter-space merging. The classic Wiener deconvolution filter reconstructs the optimal scale factor $\Gamma(u, v)$ by minimizing the mean squared error (MSE) between the calibrated spectrum and the target spectrum:
\begin{equation}
  \Gamma_{Wiener}(u, v) = \frac{1}{h(u, v)} \frac{|h(u, v)|^2}{|h(u, v)|^2 + \frac{1}{\text{SNR}(u, v)}}
\end{equation}
where $\text{SNR}(u, v)$ is the Signal-to-Noise Ratio. To simplify notation, let $M_T = \mathcal{M}_{target}(u, v)$ and $M_M = \mathcal{M}_{merged}(u, v)$ represent the target and merged spectral components at frequency $(u, v)$, respectively. Approximating the local transfer function as $h(u, v) = M_M / M_T$ and assuming a scale-invariant constant noise-to-signal ratio $c^2 = 1/\text{SNR}(u, v)$, we substitute $h(u, v)$ back into the Wiener formulation to derive the \textbf{WRSA Scaling Map}:
\begin{equation}
  \label{eq:wrsa}
  \Gamma_{WRSA}(u, v) = \frac{M_T M_M}{M_M^2 + c^2 M_T^2}
\end{equation}
Where $c > 0$ is a scale-invariant regularization parameter.

\subsection{The Stability Guarantee of WRSA}
We now prove that our scale-invariant WRSA formulation guarantees absolute numerical stability, completely bypassing the need for heuristic hard-clamping.

\begin{theorem}[WRSA Bounded Scaling Theorem]
  Let $x = \mathcal{M}_{merged}(u, v) \ge 0$ and $y = \mathcal{M}_{target}(u, v) \ge 0$. For any scale-invariant regularization parameter $c > 0$, the WRSA scaling factor is analytically bounded by:
  \begin{equation}
    \Gamma_{WRSA}(u, v) \le \frac{1}{2c}
  \end{equation}
  independent of the absolute magnitudes of the target or merged spectral components.
\end{theorem}

\begin{proof}
  See \cref{app:proof_bounded_scaling} for the complete mathematical proof.
\end{proof}

This theorem establishes a profound connection between heuristic clipping and Wiener filter theory. Setting $c = 0.1$ analytically guarantees a maximum scale amplification of exactly $\frac{1}{0.2} = 5.0$, matching FDSA's clamp threshold ($\gamma_{max} = 5.0$) but achieving this limit via a smooth, mathematically optimal roll-off. When the merged magnitude is large ($x \gg c y$), WRSA converges to the exact alignment ratio $y/x$. Conversely, when the merged magnitude collapses due to noise or sparsity ($x \to 0$), the scaling factor smoothly decays to zero:
\begin{equation}
  \lim_{x \to 0} \Gamma_{WRSA}(u, v) = 0
\end{equation}
thereby successfully suppressing high-frequency noise amplification.

\begin{theorem}[Spectral Variance Explosion and Regularization Bound]
  Let the true target spectral magnitude at frequency $(u, v)$ be $y > 0$, and let the true merged spectral magnitude be $x \ge 0$. Suppose that during calibration, the empirical merged magnitude is estimated with zero-mean additive noise $\eta$, i.e., $\hat{x} = x + \eta$, where $\eta$ is a continuous random variable with a probability density function $p(\eta)$ that is strictly positive in a neighborhood of $-x$.
  At test-time, let the merged spectral magnitude be a random variable $X_{test}$ with $\mathbb{E}[X_{test}] = x$ and $\mathbb{E}[X_{test}^2] < \infty$.
  Then, the expected mean squared reconstruction error under the naive FDSA estimator, given by:
  \begin{equation}
    \mathcal{E}(\Gamma_{FDSA}) = \mathbb{E}_{\eta, X_{test}} \left[ \left( \Gamma_{FDSA}(\hat{x}) X_{test} - y \right)^2 \right]
  \end{equation}
  is infinite ($\mathcal{E}(\Gamma_{FDSA}) = \infty$) in the sparse regime $x \to 0$. Conversely, the expected mean squared reconstruction error under the proposed WRSA estimator:
  \begin{equation}
    \mathcal{E}(\Gamma_{WRSA}) = \mathbb{E}_{\eta, X_{test}} \left[ \left( \Gamma_{WRSA}(\hat{x}) X_{test} - y \right)^2 \right]
  \end{equation}
  is strictly bounded for any scale-invariant regularization parameter $c > 0$.
\end{theorem}

\begin{proof}
  See \cref{app:proof_variance_explosion} for the complete mathematical proof.
\end{proof}

\begin{theorem}[Suboptimal Regularization Sensitivity Bound]
  Let $\mathcal{E}_{WRSA}(c^2)$ be the expected mean squared reconstruction error of the WRSA estimator under a regularization parameter $c^2 > 0$ at frequency $(u, v)$. Let $r(u, v) = P_N(u, v)/P_{target}(u, v) > 0$ be the true, frequency-dependent noise-to-signal ratio, and let $P_{target}(u, v) = \mathbb{E}[|\mathcal{M}_{target}(u, v)|^2]$ and $h(u, v)$ be the signal power and spatial degradation function respectively. Suppose the chosen regularization parameter is scaled by a factor of $\alpha > 0$ relative to the optimum, i.e., $c^2 = \alpha r(u, v)$. Then, the ratio of the suboptimal reconstruction error to the optimal reconstruction error is bounded by:
  \begin{equation}
    \frac{\mathcal{E}_{WRSA}(\alpha r)}{\mathcal{E}_{WRSA}(r)} \le \frac{(\alpha + 1)^2}{4\alpha}
  \end{equation}
  Furthermore, this upper bound is tight and is achieved when the spatial degradation function satisfies $h(u, v)^2 = \alpha r(u, v)$.
\end{theorem}

\begin{proof}
  See \cref{app:proof_sensitivity_bound} for the complete mathematical proof.
\end{proof}

This theorem provides a profound guarantee of robustness for the WRSA algorithm. It proves that even if the constant regularization parameter $c^2$ deviates substantially from the true coordinate-wise noise-to-signal ratio $r(u,v)$, the performance degradation remains extremely benign. For instance, if the parameter is off by a factor of two ($\alpha = 2$ or $\alpha = 0.5$), the mean squared reconstruction error increases by at most $12.5\%$ over the absolute theoretical optimum. This explains why a flat constant regularization parameter $c$ behaves so robustly across diverse layers and datasets.

\begin{theorem}[Spectral Estimation Lipschitz Stability Theorem]
  Let $\tilde{\Gamma}_{WRSA}(z) = \frac{z}{z^2 + c^2}$ represent the scale-normalized WRSA scaling function as a function of the merged-to-target spectral ratio $z = x/y \ge 0$, for a given scale-invariant regularization parameter $c > 0$. The sensitivity of the scaling map is globally Lipschitz continuous with a Lipschitz constant of exactly $\frac{1}{c^2}$:
  \begin{equation}
    |\tilde{\Gamma}_{WRSA}(z_1) - \tilde{\Gamma}_{WRSA}(z_2)| \le \frac{1}{c^2} |z_1 - z_2|
  \end{equation}
  for all $z_1, z_2 \ge 0$.
  In contrast, the scale-normalized naive FDSA scaling function $\tilde{\Gamma}_{FDSA}(z) = 1/z$ (for $z > \epsilon$) is not Lipschitz continuous on any interval containing $0$. Its local sensitivity grows as $\mathcal{O}(1/z^2)$, which diverges to infinity as $z \to 0$.
\end{theorem}

\begin{proof}
  See \cref{app:proof_lipschitz_stability} for the complete mathematical proof.
\end{proof}

This theorem provides an elegant and mathematically rigorous foundation for the empirical stability of WRSA. By proving that the scaling function is globally Lipschitz with constant $\frac{1}{c^2}$, we establish that any error in estimating the merged spectral profile (e.g., due to a small calibration dataset) can only propagate to the scaling factors in a strictly bounded, linear fashion. In contrast, naive spectral methods suffer from unbounded sensitivity in sparse frequency bands, mathematically explaining why FDSA is highly unstable across random calibration seeds.

\begin{theorem}[Phase Coherence and Spatial Edge Preservation Theorem]
  Let $X \in \mathbb{R}^{H \times W}$ be a spatial activation map with 2D Discrete Fourier Transform $\tilde{X}(u, v) = A(u, v) e^{i \phi(u, v)}$. Let the target expert activation map have spectrum $\tilde{Y}(u, v) = B(u, v) e^{i \theta(u, v)}$. Under a phase-preserving spectral calibration operator $\mathcal{C}$ that rescales magnitudes by $\Gamma(u, v) \ge 0$ while leaving phase intact (i.e., $\tilde{X}_{calibrated}(u, v) = \Gamma(u, v) A(u, v) e^{i \phi(u, v)}$), the spatial $L_2$ reconstruction error is minimized if and only if the phase of the merged representation is fully coherent with the target (i.e., $\phi(u, v) = \theta(u, v)$). Furthermore, any phase perturbation $\Delta \phi(u, v)$ shifts and distorts spatial activation boundaries, leading to representational mismatch in subsequent non-linear layers.
\end{theorem}

\begin{proof}
  See \cref{app:proof_phase_coherence} for the complete mathematical proof.
\end{proof}

This theorem provides a profound mathematical justification for a fundamental design choice shared by both WRSA and FDSA: performing representation calibration exclusively on the Fourier spectral magnitudes while keeping the empirical phase of the merged activations completely unmodified. Because the task experts are fine-tuned from a shared, pre-trained progenitor network, they retain highly aligned spatial phase structures. Modifying the phase would destroy this spatial edge coherence, whereas WRSA's phase-preserving magnitude scaling stably restores collapsed high-frequency features without introducing geometric phase distortion.

\section{Experimental Evaluation}
\label{sec:experiments}

\subsection{Experimental Setup}
We evaluate our method on the standard multi-task vision benchmark consisting of MNIST, Fashion-MNIST, and CIFAR-10.
\textbf{Model \& Training}: We utilize a pre-trained ResNet-18 backbone. Three specialized task expert models are trained sequentially on deterministic subsets of 3,000 images from each dataset for 5 epochs using the AdamW optimizer with learning rate $5 \times 10^{-4}$ and weight decay $10^{-4}$.
\textbf{Calibration Protocol}: We construct a joint calibration dataset of size $3 \times 128 = 384$ samples by pooling 128 samples per task.
\textbf{Evaluation Metric}: Accuracy is reported on the full 10,000-sample test sets of each dataset. Evaluating a merged model on task $t$ is performed by combining the merged backbone with the task-specific expert head.

\subsection{Quantitative Merging Results}
The quantitative results for Weight Averaging (WA) and Task Arithmetic (TA) under different calibration methods are detailed in \cref{tab:results}.

\begin{table*}[t]
\caption{Multi-task classification accuracy (\%) on the ResNet-18 benchmark under Weight Averaging (WA) and Task Arithmetic (TA, $\lambda = 0.3$) merging. We compare our proposed Wiener-Regularized Spectral Alignment (WRSA) against uncalibrated, spatial-calibrated (SP-TAAC) \cite{submission9}, and naive spectral-calibrated (FDSA) \cite{submission10} methods.}
\label{tab:results}
\vskip 0.15in
\begin{center}
\begin{small}
\begin{tabular}{llcccc}
\toprule
\textbf{Merge Mode} & \textbf{Calibration Method} & \textbf{MNIST} & \textbf{F-MNIST} & \textbf{CIFAR-10} & \textbf{Average} \\
\midrule
\textbf{Oracle} & Task Experts (Upper Bound) & 97.21 & 84.02 & 64.36 & 81.86 \\
\midrule
\textbf{WA} & Uncalibrated Baseline & 58.56 & 54.43 & 31.18 & 48.06 \\
\textbf{WA} & SP-TAAC \cite{submission9} (Spatial) & 60.09 & 67.85 & 34.57 & \textbf{54.17} \\
\textbf{WA} & FDSA \cite{submission10} (Naive Spectral) & 29.93 & 61.94 & 34.10 & 41.99 \\
\textbf{WA} & \textbf{WRSA (Ours, $c=0.01$)} & 29.98 & 61.95 & 34.13 & 42.02 \\
\textbf{WA} & \textbf{WRSA (Ours, $c=0.05$)} & 31.32 & 62.74 & 34.44 & 42.83 \\
\textbf{WA} & \textbf{WRSA (Ours, $c=0.10$)} & 35.80 & 64.51 & 35.56 & 45.29 \\
\textbf{WA} & \textbf{WRSA (Ours, $c=0.20$)} & 48.61 & 68.61 & 37.98 & 51.73 \\
\textbf{WA} & \textbf{WRSA (Ours, $c=0.30$)} & 54.55 & 70.06 & 37.77 & \textbf{54.13} \\
\textbf{WA} & \textbf{WRSA (Ours, $c=0.40$)} & 55.39 & 57.03 & 26.21 & 46.21 \\
\textbf{WA} & \textbf{WRSA (Ours, $c=0.50$)} & 30.37 & 20.14 & 16.35 & 22.29 \\
\midrule
\textbf{TA} & Uncalibrated Baseline & 54.38 & 56.29 & 32.26 & 47.64 \\
\textbf{TA} & SP-TAAC \cite{submission9} (Spatial) & 59.04 & 67.33 & 36.37 & \textbf{54.25} \\
\textbf{TA} & FDSA \cite{submission10} (Naive Spectral) & 37.68 & 52.24 & 36.26 & 42.06 \\
\textbf{TA} & \textbf{WRSA (Ours, $c=0.01$)} & 37.69 & 52.28 & 36.28 & 42.08 \\
\textbf{TA} & \textbf{WRSA (Ours, $c=0.05$)} & 38.67 & 53.09 & 36.64 & 42.80 \\
\textbf{TA} & \textbf{WRSA (Ours, $c=0.10$)} & 41.99 & 55.65 & 38.01 & 45.22 \\
\textbf{TA} & \textbf{WRSA (Ours, $c=0.20$)} & 50.20 & 63.53 & 40.25 & \textbf{51.33} \\
\bottomrule
\end{tabular}
\end{small}
\end{center}
\vskip -0.1in
\end{table*}

\subsection{Ablation and Stability Studies}
To rigorously evaluate the sample complexity and robustness of our proposed framework in alignment with our theoretical claims, we conduct two extensive studies: (i) a sample size ablation study where the size of the calibration dataset $N$ per task is varied across $\{16, 64, 128\}$ (resulting in total calibration sets of $\{48, 192, 384\}$ samples), and (ii) a multi-seed statistical stability study over 5 different randomly sampled calibration runs.

\begin{table}[tb]
\caption{Sample size ablation study under Weight Averaging (WA) merging. We evaluate multi-task accuracy (\%) for calibration sizes $N \in \{16, 64, 128\}$ per task.}
\label{tab:ablation}
\vskip 0.1in
\centering
\begin{scriptsize}
\setlength{\tabcolsep}{2pt}
\begin{tabular}{llcccc}
\toprule
\textbf{N} & \textbf{Method} & \textbf{MNIST} & \textbf{F-MNIST} & \textbf{CIFAR-10} & \textbf{Average} \\
\midrule
16 & Uncalibrated & 58.56 & 54.43 & 31.18 & 48.06 \\
16 & SP-TAAC & 58.52 & 67.97 & 34.84 & 53.78 \\
16 & FDSA & 41.62 & 61.44 & 36.08 & 46.38 \\
16 & \textbf{WRSA (Ours)} & \textbf{48.63} & \textbf{67.67} & \textbf{38.27} & \textbf{51.52} \\
\midrule
64 & Uncalibrated & 58.56 & 54.43 & 31.18 & 48.06 \\
64 & SP-TAAC & 59.40 & 68.19 & 34.56 & 54.05 \\
64 & FDSA & 42.64 & 62.00 & 36.46 & 47.03 \\
64 & \textbf{WRSA (Ours)} & \textbf{48.11} & \textbf{67.66} & \textbf{38.07} & \textbf{51.28} \\
\midrule
128 & Uncalibrated & 58.56 & 54.43 & 31.18 & 48.06 \\
128 & SP-TAAC & 58.54 & 67.87 & 34.17 & 53.53 \\
128 & FDSA & 42.32 & 61.80 & 36.21 & 46.78 \\
128 & \textbf{WRSA (Ours)} & \textbf{48.83} & \textbf{67.66} & \textbf{37.86} & \textbf{51.45} \\
\bottomrule
\end{tabular}
\end{scriptsize}
\vskip -0.1in
\end{table}

\begin{table*}[t]
\caption{Multi-seed statistical stability and variance analysis under Weight Averaging (WA) with $N=128$. We compare standard FDSA with our proposed WRSA ($c=0.30$) across 5 random calibration seeds.}
\label{tab:stability}
\vskip 0.1in
\begin{center}
\begin{small}
\begin{tabular}{lcccccccc}
\toprule
\textbf{Method} & \textbf{Seed 42} & \textbf{Seed 43} & \textbf{Seed 44} & \textbf{Seed 45} & \textbf{Seed 46} & \textbf{Mean Acc} & \textbf{Std Dev} & \textbf{Var. Reduction} \\
\midrule
FDSA & 46.08 & 46.30 & 47.28 & 45.59 & 45.41 & 46.13 & 0.6577 & Baseline \\
\textbf{WRSA (Ours)} & \textbf{51.75} & \textbf{51.72} & \textbf{51.74} & \textbf{52.11} & \textbf{52.14} & \textbf{51.89} & \textbf{0.1924} & \textbf{3.4$\times$ Lower} \\
\bottomrule
\end{tabular}
\end{small}
\end{center}
\vskip -0.1in
\end{table*}

As shown in \cref{tab:ablation}, FDSA's performance collapses under small calibration sizes, yielding only 46.38\% average accuracy at $N=16$ (a 2.18\% drop compared to $N=128$). Under extreme sparsity ($N=16$), pointwise division in FDSA suffers from high estimation noise, particularly on MNIST where accuracy plummets to 41.62\%. In contrast, WRSA maintains stable performance. At $N=16$, WRSA achieves \textbf{51.52\%} average accuracy, outperforming FDSA by \textbf{+5.14\%} absolute and demonstrating that Wiener regularization successfully insulates representation alignment from small-data estimation noise.

Furthermore, the statistical stability results in \cref{tab:stability} reveal that WRSA dramatically reduces estimation variance across different random calibration subsets. Across 5 random seeds, FDSA exhibits a high standard deviation of \textbf{0.6577}, with its average accuracy fluctuating between 45.41\% and 47.28\%. In contrast, WRSA achieves a mean accuracy of \textbf{51.89\%} with a remarkably low standard deviation of only \textbf{0.1924}---a \textbf{3.4$\times$ reduction in variance} over FDSA. Looking at individual tasks, WRSA reduces the standard deviation on MNIST from 1.2010\% (FDSA) down to \textbf{0.6454\%}, and on Fashion-MNIST from 0.4868\% down to \textbf{0.0893\%} (a \textbf{5.4$\times$ reduction} in standard deviation). This provides definitive empirical verification that our scale-invariant Wiener filter acts as a mathematically sound regularizer, guaranteeing statistical stability and reliable representation recovery under any choice of calibration subset.

\section{Discussion and Results Analysis}
\label{sec:discussion}

\subsection{Deconstructing the Spectral Sparsity Trap}
As shown in \cref{tab:results}, standard FDSA suffers a catastrophic collapse on the MNIST dataset under Weight Averaging, falling from an uncalibrated accuracy of \textbf{58.56\%} down to \textbf{29.93\%}. This empirical result provides clear proof of the \emph{Spectral Sparsity Trap} hypothesized in \cref{sec:theory}. Because MNIST images are extremely sparse (comprising localized bright digits on flat black backdrops), their activation maps are mostly zero. 
Pointwise magnitude division in FDSA divides by these collapsed merged magnitudes, heavily amplifying estimating noise and scrambling the spatial structure of deeper feature maps.

\begin{figure}[htbp]
  \centering
  \includegraphics[width=0.88\linewidth]{plots/wrsa_tradeoff.pdf}
  \caption{Empirical trade-off curve of WRSA Average Multi-Task Accuracy vs. the regularization parameter $c$, compared to Uncalibrated, SP-TAAC, and FDSA baselines under Weight Averaging.}
  \label{fig:tradeoff}
\end{figure}

\subsection{The WRSA Regularization Trade-off Curve}
By sweeping $c$ in \cref{fig:tradeoff}, we identify three distinct regimes of spectral alignment behavior:
(1)~\textbf{Noise Amplification Regime ($c \le 0.05$)}: For very small values of $c$, the Wiener filter approaches naive pointwise division. High-frequency noise is amplified, leading to poor performance on MNIST and depressed overall averages (~42\%).
(2)~\textbf{Optimal Signal Restoration ($c \in [0.20, 0.30]$)}: As $c$ increases, WRSA smoothly regularizes the scaling profile. At $c=0.30$, WRSA achieves a remarkable average accuracy of \textbf{54.13\%} under WA merging, which represents a massive \textbf{+12.14\% absolute improvement} over standard FDSA (41.99\%) and matches the spatial-domain SP-TAAC (54.17\%) while operating entirely in the frequency domain. Notably, on Fashion-MNIST, WRSA ($c=0.30$) achieves \textbf{70.06\%}, outperforming even SP-TAAC (\textbf{67.85\%}).
(3)~\textbf{Signal Collapse Regime ($c \ge 0.40$)}: When $c$ is set too high, the regularization term in \cref{eq:wrsa} dominates, causing the scaling factors to shrink globally towards zero. This attenuates activation maps to zero, causing the network's deep features to collapse and performance to plummet to random guessing (10.03\%).

These three regimes show that frequency-domain model merging calibration requires a careful, mathematically sound balance between high-frequency signal restoration and noise suppression. WRSA provides the exact analytical knob to navigate this trade-off optimally.

\subsection{Active Inference Overhead and Parameter Fusion Limitations}
An important aspect of representation calibration is its deployability in resource-constrained environments. While spatial-domain calibration methods like SP-TAAC \cite{submission9} or REPAIR \cite{jordan2022repair} perform global or channel-wise affine transformations, these scaling factors are static and can be fused directly back into the preceding layer's weights and biases (e.g., $W' = \gamma W$, $b' = \gamma b$). Consequently, spatial calibration achieves \emph{Zero Inference Overhead (ZIO)} at test-time, leaving the computational cost of the forward pass unaltered.

In contrast, spectral-domain methods such as WRSA and FDSA apply a frequency-dependent scaling map $\Gamma(u, v)$ directly to the Fourier coefficients of intermediate representations. This operation requires a 2D Discrete Fourier Transform (FFT) and its inverse (IFFT) along the spatial dimensions at each calibrated layer during active inference. For a feature map $X \in \mathbb{R}^{B \times C \times H \times W}$, this adds a computational complexity of $\mathcal{O}(B \cdot C \cdot H W \log(H W))$ floating-point operations (FLOPs).

Furthermore, we prove that this active inference overhead is theoretically impossible to eliminate via spatial parameter fusion.
\begin{theorem}[Spectral Non-Fusion Theorem]
  Let $\Gamma(u, v) \in \mathbb{R}^{H \times W}$ be a non-flat frequency-domain scaling map, and let $*$ denote circular spatial convolution. The spectral calibration operation $\mathcal{F}^{-1}(\Gamma \cdot \mathcal{F}(X))$ is equivalent to a spatial circular convolution of $X$ with a filter $G = \mathcal{F}^{-1}(\Gamma) \in \mathbb{R}^{H \times W}$. If $\Gamma$ is non-flat, $G$ is not a spatially localized delta impulse (i.e., its spatial support is the entire $H \times W$ grid). Therefore, the operation cannot be represented as a localized spatial convolution (such as a $3 \times 3$ or $1 \times 1$ kernel), making it impossible to fuse $\Gamma$ into standard localized convolutional weights.
\end{theorem}
\begin{proof}
  See \cref{app:proof_non_fusion} for the complete mathematical proof.
\end{proof}
This theorem formalizes a fundamental trade-off in representation calibration: while frequency-domain methods like WRSA offer superior precision and granular control over collapsed representational structures, they require paying an active computational overhead of $\mathcal{O}(B \cdot C \cdot H W \log(H W))$ at test-time due to the mathematical impossibility of spatial parameter fusion.

To empirically quantify this active inference overhead, we benchmark the forward pass latency of the ResNet-18 backbone on CPU using a batch size of $128$ over $100$ timed iterations (following $50$ warmup steps). The empirical latencies and relative overheads are reported in \cref{tab:latency}. While SP-TAAC achieves Zero Inference Overhead (ZIO) through in-place parameter fusion, the active 2D FFT/IFFT operations of FDSA and WRSA across all $20$ BatchNorm layers of the ResNet-18 backbone increase the latency from $45.97$\,ms to $93.37$\,ms ($2.03\times$) and $94.39$\,ms ($2.05\times$) respectively. This $2\times$ latency penalty is extremely consistent with our theoretical $\mathcal{O}(B \cdot C \cdot H W \log(H W))$ complexity analysis, demonstrating that while spectral methods yield superior accuracy and noise insulation, they introduce a distinct computational trade-off.
\begin{table}[tb]
\caption{Inference latency and relative computational overhead of different calibration methods on ResNet-18 under a batch size of $128$ on CPU. We compare our proposed WRSA against uncalibrated, SP-TAAC \cite{submission9}, and FDSA \cite{submission10} methods.}
\label{tab:latency}
\vskip 0.1in
\centering
\begin{scriptsize}
\setlength{\tabcolsep}{2pt}
\begin{tabular}{lcccc}
\toprule
\textbf{Method} & \textbf{Hooks} & \textbf{Fusion} & \textbf{Lat. (ms)} & \textbf{Overhead} \\
\midrule
Uncalibrated & No & N/A & 45.97 & $1.00\times$ \\
SP-TAAC & No & Yes & 45.97 & $1.00\times$ (ZIO) \\
FDSA & Yes & No & 93.37 & $2.03\times$ \\
\textbf{WRSA (Ours)} & Yes & No & 94.39 & $2.05\times$ \\
\bottomrule
\end{tabular}
\end{scriptsize}
\vskip -0.1in
\end{table}

\subsection{Theoretical Assumptions and Methodological Limitations}
While WRSA provides a mathematically sound and highly stable framework for spectral representation alignment, its design relies on several key theoretical assumptions and simplifications that present opportunities for future research:
(i)~\textbf{Scale-Invariant Noise Approximation}: Our derivation in \cref{eq:wrsa} approximates the noise-to-signal ratio $P_N(u, v)/P_{target}(u, v)$ as a flat constant $c^2$. In reality, the estimation noise from a small calibration dataset may exhibit non-stationary behavior across frequency bands. Future work could dynamically estimate $P_N(u, v)$ across coordinates using empirical variance.
(ii)~\textbf{Spatial Translation Invariance}: The 2D DFT assumes translation-invariant spatial structures, natural for CNNs. However, in non-grid architectures (e.g., Graph Neural Networks or fully connected layers), this breaks down, requiring alternative spectral bases like Graph Fourier Transforms or SVD.
(iii)~\textbf{Active Inference Overhead}: As proven by the Spectral Non-Fusion Theorem, frequency-dependent scaling maps cannot be folded into standard localized spatial parameters, creating an active inference overhead of $\mathcal{O}(B \cdot C \cdot H W \log(H W))$ per layer. For real-time applications, this presents a trade-off between calibration accuracy and throughput.
Evaluating these boundary conditions ensures a balanced assessment of spectral calibration, guiding the field towards more generalizable merging frameworks.

\section{Conclusion \& Future Work}
\label{sec:conclusion}
In this work, we deconstructed frequency-domain activation calibration in model merging and identified a critical failure mode: the \emph{Spectral Sparsity Trap}. We introduced \textbf{Wiener-Regularized Spectral Alignment (WRSA)}, which resolves the noise amplification bottleneck of naive pointwise division. We proved that our formulation analytically bounds the spectral scaling factors, replacing heuristic hard-clamping with a smooth, mathematically optimal Wiener roll-off. Extensive empirical evaluation demonstrated that WRSA outclasses naive spectral methods by up to \textbf{+12.14\% absolute accuracy}, proving that frequency-domain calibration can be both theoretically sound and empirically highly competitive. 

Future work will focus on extending the Wiener deconvolution framework to Transformer architectures, where attention matrices and token sequences suffer from spectral degradation during Large Language Model (LLM) merging.

\nocite{*}
\bibliography{paper}
\bibliographystyle{icml2026}

\newpage
\appendix
\onecolumn

\section{Formal Derivation of the WRSA Wiener Filter}
\label{app:wiener_derivation}
In this section, we provide the full, rigorous mathematical derivation of the Wiener spectral scaling map $\Gamma_{WRSA}(u, v)$ from first principles.
We model the observed merged model's activation spectral components in the 2D Fourier domain as a degraded version of the target (optimal expert) spectral components, subject to additive random estimation noise:
\begin{equation}
  \mathcal{M}_{merged}(u, v) = h(u, v) \mathcal{M}_{target}(u, v) + N(u, v)
\end{equation}
where $h(u, v) \in [0, 1]$ represents the linear spatial low-pass transfer function induced by parameter interference, and $N(u, v)$ is a zero-mean random noise process with variance $\sigma_N^2(u, v)$ representing the statistical error of estimating spectral density from a finite calibration dataset of size $N$.

We seek a linear spectral scaling map $\Gamma(u, v)$ to reconstruct the target spectrum:
\begin{equation}
  \hat{\mathcal{M}}_{target}(u, v) = \Gamma(u, v) \mathcal{M}_{merged}(u, v)
\end{equation}
such that the expected mean squared error (MSE) of the reconstruction is minimized:
\begin{equation}
  E[\Gamma(u, v)] = \mathbb{E} \left[ \left| \hat{\mathcal{M}}_{target}(u, v) - \mathcal{M}_{target}(u, v) \right|^2 \right]
\end{equation}
Expanding this quadratic expectation:
\begin{align*}
  E[\Gamma] &= \mathbb{E} \left[ \left| \Gamma (h \mathcal{M}_{target} + N) - \mathcal{M}_{target} \right|^2 \right] \\
  &= \mathbb{E} \left[ \left| (\Gamma h - 1) \mathcal{M}_{target} + \Gamma N \right|^2 \right]
\end{align*}
Under the standard assumption that the signal $\mathcal{M}_{target}$ and the estimation noise $N$ are statistically orthogonal ($\mathbb{E}[\mathcal{M}_{target} N^*] = 0$), the cross-terms vanish:
\begin{align*}
  E[\Gamma] &= (\Gamma h - 1)^2 \mathbb{E}[|\mathcal{M}_{target}|^2] + \Gamma^2 \mathbb{E}[|N|^2] \\
  &= (\Gamma h - 1)^2 P_{target} + \Gamma^2 P_N
\end{align*}
where $P_{target}(u, v) = \mathbb{E}[|\mathcal{M}_{target}(u, v)|^2]$ is the power spectral density of the target activations, and $P_N(u, v) = \sigma_N^2(u, v)$ is the noise power spectral density.

To find the optimal filter $\Gamma^*(u, v)$ that minimizes $E[\Gamma]$, we take the partial derivative with respect to $\Gamma$ and set it to zero:
\begin{equation}
  \frac{\partial E[\Gamma]}{\partial \Gamma} = 2 h (\Gamma h - 1) P_{target} + 2 \Gamma P_N = 0
\end{equation}
Solving for $\Gamma$:
\begin{align*}
  \Gamma (h^2 P_{target} + P_N) &= h P_{target} \\
  \Gamma^*(u, v) &= \frac{h(u, v) P_{target}(u, v)}{h(u, v)^2 P_{target}(u, v) + P_N(u, v)} \\
  &= \frac{1}{h(u, v)} \frac{h(u, v)^2}{h(u, v)^2 + \frac{P_N(u, v)}{P_{target}(u, v)}}
\end{align*}
This is the classic formulation of the Wiener deconvolution filter. 

To make this formulation scale-invariant and data-driven for model merging, we approximate the local degradation function as the pointwise ratio:
\begin{equation}
  h(u, v) \approx \frac{\mathcal{M}_{merged}(u, v)}{\mathcal{M}_{target}(u, v)}
\end{equation}
and model the noise-to-signal ratio as a scale-invariant constant:
\begin{equation}
  \frac{P_N(u, v)}{P_{target}(u, v)} \approx c^2
\end{equation}
where $c > 0$ represents the inverse of the Signal-to-Noise Ratio (SNR). Substituting these approximations back into the Wiener filter:
\begin{align*}
  \Gamma^*(u, v) &= \frac{\mathcal{M}_{target}(u, v)}{\mathcal{M}_{merged}(u, v)} \frac{\left( \frac{\mathcal{M}_{merged}(u, v)}{\mathcal{M}_{target}(u, v)} \right)^2}{\left( \frac{\mathcal{M}_{merged}(u, v)}{\mathcal{M}_{target}(u, v)} \right)^2 + c^2} \\
  &= \frac{\mathcal{M}_{target}(u, v)}{\mathcal{M}_{merged}(u, v)} \frac{\mathcal{M}_{merged}(u, v)^2}{\mathcal{M}_{merged}(u, v)^2 + c^2 \mathcal{M}_{target}(u, v)^2} \\
  &= \frac{\mathcal{M}_{target}(u, v) \mathcal{M}_{merged}(u, v)}{\mathcal{M}_{merged}(u, v)^2 + c^2 \mathcal{M}_{target}(u, v)^2}
\end{align*}
This completes the formal proof of the Wiener-Regularized Spectral Alignment (WRSA) formula.

\section{Theoretical Link between Spatial and Spectral Calibration}
\label{app:parseval_link}
We establish a theoretical link between spatial-domain variance alignment (like SP-TAAC) and frequency-domain spectral alignment (WRSA) using Parseval's Theorem.
Let $X \in \mathbb{R}^{H \times W}$ be a zero-mean intermediate activation map, and let $\tilde{X} \in \mathbb{C}^{H \times W}$ be its 2D Fourier transform. Parseval's theorem asserts that the total energy is conserved across domains:
\begin{equation}
  \sum_{y=1}^H \sum_{x=1}^W |X(y, x)|^2 = \frac{1}{H \cdot W} \sum_{v=1}^H \sum_{u=1}^W |\tilde{X}(v, u)|^2
\end{equation}
Because $X$ has zero mean, the empirical spatial variance $\sigma^2$ is directly proportional to its total energy:
\begin{equation}
  \sigma^2 = \frac{1}{H^2 \cdot W^2} \sum_{v=1}^H \sum_{u=1}^W |\tilde{X}(v, u)|^2
\end{equation}
Under a spatial-domain calibration method (e.g. SP-TAAC), the entire activation map is multiplied by a flat global scalar $\gamma_{spatial} = \sigma_{target} / \sigma_{merged}$. In the frequency domain, this is equivalent to scaling all Fourier coefficients by the same flat scalar:
\begin{equation}
  \tilde{X}^*_{spatial}(v, u) = \gamma_{spatial} \cdot \tilde{X}(v, u)
\end{equation}
In contrast, our proposed WRSA applies a frequency-dependent scaling map $\Gamma_{WRSA}(v, u)$ to each coordinate of the 2D Fourier transform:
\begin{equation}
  \tilde{X}^*_{WRSA}(v, u) = \Gamma_{WRSA}(v, u) \cdot \tilde{X}(v, u)
\end{equation}
This reveals that spatial-domain calibration is a special, restricted case of spectral alignment where the spectral scaling map is restricted to be a flat constant, i.e., $\Gamma(v, u) = \gamma_{spatial}$ for all frequencies. 
By resolving the scaling factor at each 2D spatial frequency $(v, u)$, WRSA generalizes spatial-domain alignment. It allows the model to selectively restore high-frequency representations that are collapsed by model merging while smoothly attenuating those corrupted by sparsity noise, a degree of freedom that is impossible under flat spatial scaling.

\section{Detailed Stability and Robustness Proofs}
\label{app:proofs_stability_robustness}

\subsection{Proof of Theorem 3.2 (WRSA Bounded Scaling Theorem)}
\label{app:proof_bounded_scaling}
\begin{proof}
  Let $f(x) = \frac{y x}{x^2 + c^2 y^2}$ for $x \ge 0$. To find the maximum of $f(x)$, we take its derivative with respect to $x$:
  \begin{equation}
    f'(x) = \frac{y(x^2 + c^2 y^2) - y x (2x)}{(x^2 + c^2 y^2)^2} = \frac{y (c^2 y^2 - x^2)}{(x^2 + c^2 y^2)^2}
  \end{equation}
  Setting the derivative to zero yields:
  \begin{equation}
    c^2 y^2 - x^2 = 0 \implies x = c y
  \end{equation}
  Substituting $x = c y$ back into $f(x)$ yields the maximum value:
  \begin{equation}
    f(c y) = \frac{y (c y)}{(c y)^2 + c^2 y^2} = \frac{c y^2}{2 c^2 y^2} = \frac{1}{2c}
  \end{equation}
  Since $f(0) = 0$ and $\lim_{x\to\infty} f(x) = 0$, the global maximum of $\Gamma_{WRSA}(u, v)$ is exactly $\frac{1}{2c}$. This completes the proof.
\end{proof}

\subsection{Proof of Theorem 3.3 (Spectral Variance Explosion and Regularization Bound)}
\label{app:proof_variance_explosion}
\begin{proof}
  First, we analyze the expected squared value of the FDSA scaling factor $\Gamma_{FDSA}(\hat{x}) = \frac{y}{\hat{x}} = \frac{y}{x + \eta}$. The expected squared value is given by:
  \begin{equation}
    \mathbb{E}_{\eta} \left[ \Gamma_{FDSA}(\hat{x})^2 \right] = y^2 \int_{-\infty}^{\infty} \frac{1}{(x + \eta)^2} p(\eta) \, d\eta
  \end{equation}
  Since $p(\eta)$ is continuous and strictly positive in a neighborhood of $-x$ (say, $[ -x - \delta, -x + \delta ]$ for some $\delta > 0$), we can write $p(\eta) \ge p_0 > 0$ for all $\eta \in [ -x - \delta, -x + \delta ]$. Thus, the integral is lower-bounded by:
  \begin{equation}
    \mathbb{E}_{\eta} \left[ \Gamma_{FDSA}(\hat{x})^2 \right] \ge y^2 p_0 \int_{-x-\delta}^{-x+\delta} \frac{1}{(x + \eta)^2} \, d\eta
  \end{equation}
  Letting $u = x + \eta$, the integral becomes:
  \begin{equation}
    \int_{-\delta}^{\delta} \frac{1}{u^2} \, du = \left[ -\frac{1}{u} \right]_{-\delta}^{\delta} = \infty
  \end{equation}
  Since the expected second moment of the scaling factor is infinite, and $X_{test}$ has finite, non-zero variance independent of $\eta$, the expected reconstruction error $\mathcal{E}(\Gamma_{FDSA})$ diverges to infinity.

  For the WRSA estimator, we invoke the WRSA Bounded Scaling Theorem (Theorem 3.2), which guarantees that for any $\hat{x}$ and $y$, we have $\Gamma_{WRSA}(\hat{x}) \le \frac{1}{2c}$. Substituting this bound into the reconstruction error expression:
  \begin{align*}
    \mathcal{E}(\Gamma_{WRSA}) &= \mathbb{E}_{\eta, X_{test}} \left[ \left( \Gamma_{WRSA}(\hat{x}) X_{test} - y \right)^2 \right] \\
    &\le \mathbb{E}_{\eta, X_{test}} \left[ 2 \Gamma_{WRSA}(\hat{x})^2 X_{test}^2 + 2 y^2 \right] \\
    &\le \mathbb{E}_{\eta, X_{test}} \left[ 2 \left( \frac{1}{2c} \right)^2 X_{test}^2 + 2 y^2 \right] \\
    &= \frac{1}{2c^2} \mathbb{E}[X_{test}^2] + 2 y^2 < \infty
  \end{align*}
  Since $\mathbb{E}[X_{test}^2]$ and $y^2$ are finite, the WRSA expected reconstruction error is strictly bounded. This mathematically guarantees the stability and robust noise-insulating capability of WRSA over naive pointwise division.
\end{proof}

\subsection{Proof of Theorem 3.4 (Suboptimal Regularization Sensitivity Bound)}
\label{app:proof_sensitivity_bound}
\begin{proof}
  Under zero-mean additive estimation noise and orthogonal signal assumptions, the expected mean squared reconstruction error for a chosen regularization parameter $c^2 > 0$ is given by:
  \begin{equation}
    \mathcal{E}_{WRSA}(c^2) = P_{target} \left[ \frac{c^4 + h^2 r}{(h^2 + c^2)^2} \right]
  \end{equation}
  Setting $c^2 = r$ yields the theoretically optimal Wiener filter reconstruction error:
  \begin{equation}
    \mathcal{E}_{WRSA}(r) = P_{target} \left[ \frac{r^2 + h^2 r}{(h^2 + r)^2} \right] = P_{target} \frac{r}{h^2 + r}
  \end{equation}
  For a suboptimal regularization $c^2 = \alpha r$, the reconstruction error is:
  \begin{equation}
    \mathcal{E}_{WRSA}(\alpha r) = P_{target} \left[ \frac{\alpha^2 r^2 + h^2 r}{(h^2 + \alpha r)^2} \right]
  \end{equation}
  Taking the ratio of the suboptimal error to the optimal error:
  \begin{align}
    R(\alpha) &= \frac{\mathcal{E}_{WRSA}(\alpha r)}{\mathcal{E}_{WRSA}(r)} \nonumber \\
    &= \frac{(\alpha^2 r^2 + h^2 r)(h^2 + r)}{r (h^2 + \alpha r)^2} \nonumber \\
    &= \frac{(\alpha^2 r + h^2)(h^2 + r)}{(h^2 + \alpha r)^2}
  \end{align}
  Let $u = h^2 / r \ge 0$. We express $R(\alpha)$ as a function of $u$:
  \begin{align}
    R(\alpha) &= \frac{(\alpha^2 + u)(u + 1)}{(u + \alpha)^2} \nonumber \\
    &= \frac{u^2 + (\alpha^2 + 1) u + \alpha^2}{(u + \alpha)^2} \nonumber \\
    &= 1 + \frac{(\alpha - 1)^2 u}{(u + \alpha)^2}
  \end{align}
  To find the maximum of $R(\alpha)$ with respect to $u \ge 0$, we analyze the function $k(u) = \frac{u}{(u + \alpha)^2}$. Taking the derivative with respect to $u$:
  \begin{equation}
    k'(u) = \frac{(u + \alpha)^2 - 2u(u + \alpha)}{(u + \alpha)^4} = \frac{\alpha - u}{(u + \alpha)^3}
  \end{equation}
  Setting $k'(u) = 0$ yields the unique global maximum at $u = \alpha$, which corresponds to $h^2 = \alpha r$. Substituting $u = \alpha$ back into our ratio expression:
  \begin{equation}
    R(\alpha) \le 1 + \frac{(\alpha - 1)^2 \alpha}{4 \alpha^2} = 1 + \frac{(\alpha - 1)^2}{4\alpha} = \frac{(\alpha + 1)^2}{4\alpha}
  \end{equation}
  This completes the proof.
\end{proof}

\subsection{Proof of Theorem 3.5 (Spectral Estimation Lipschitz Stability Theorem)}
\label{app:proof_lipschitz_stability}
\begin{proof}
  We first analyze the derivative of the scale-normalized WRSA scaling function $\tilde{\Gamma}_{WRSA}(z)$ with respect to $z$:
  \begin{align}
    \tilde{\Gamma}'_{WRSA}(z) &= \frac{d}{dz}\left( \frac{z}{z^2 + c^2} \right) \nonumber \\
    &= \frac{c^2 - z^2}{(z^2 + c^2)^2}
  \end{align}
  By the Mean Value Theorem, bounding the Lipschitz constant of our scale-normalized map by $L = 1/c^2$ on the domain $[0, \infty)$ is equivalent to showing:
  \begin{equation}
    \sup_{z \ge 0} \left| \tilde{\Gamma}'_{WRSA}(z) \right| = \frac{1}{c^2}
  \end{equation}
  We analyze the function $g(z) = \tilde{\Gamma}'_{WRSA}(z) = \frac{c^2 - z^2}{(z^2 + c^2)^2}$.
  Since $g(z)$ is continuous and differentiable on $[0, \infty)$, we identify its critical points by taking its derivative with respect to $z$:
  \begin{align*}
    g'(z) &= \frac{-2z(z^2 + c^2)^2 - (c^2 - z^2) \cdot 2(z^2 + c^2) \cdot 2z}{(z^2 + c^2)^4} \\
    &= \frac{-2z(z^2 + c^2) - 4z(c^2 - z^2)}{(z^2 + c^2)^3} \\
    &= \frac{-2z^3 - 2c^2 z - 4c^2 z + 4z^3}{(z^2 + c^2)^3} \\
    &= \frac{2z(z^2 - 3c^2)}{(z^2 + c^2)^3}
  \end{align*}
  Setting $g'(z) = 0$ for $z \ge 0$ yields two critical points: $z = 0$ and $z = \sqrt{3}c$.
  
  Now, we evaluate the absolute value of $g(z)$ at the boundaries and critical points:
  \begin{enumerate}
    \item At $z = 0$: 
    \begin{equation}
      |g(0)| = \left| \frac{c^2 - 0}{(0 + c^2)^2} \right| = \frac{c^2}{c^4} = \frac{1}{c^2}
    \end{equation}
    \item At $z = \sqrt{3}c$:
    \begin{equation}
      |g(\sqrt{3}c)| = \left| \frac{c^2 - 3c^2}{(3c^2 + c^2)^2} \right| = \left| \frac{-2c^2}{16c^4} \right| = \frac{1}{8c^2}
    \end{equation}
    \item At the limit as $z \to \infty$:
    \begin{equation}
      \lim_{z \to \infty} |g(z)| = 0
    \end{equation}
  \end{enumerate}
  Comparing these values, the maximum absolute value is achieved at $z = 0$, giving:
  \begin{equation}
    \sup_{z \ge 0} \left| \tilde{\Gamma}'_{WRSA}(z) \right| = \frac{1}{c^2}
  \end{equation}
  Therefore, $\tilde{\Gamma}_{WRSA}(z)$ is globally Lipschitz continuous on $[0, \infty)$ with Lipschitz constant $L = \frac{1}{c^2}$.
  
  For naive FDSA, the scale-normalized scaling map is $\tilde{\Gamma}_{FDSA}(z) = \frac{1}{z}$. Its derivative with respect to $z$ is:
  \begin{equation}
    \tilde{\Gamma}'_{FDSA}(z) = -\frac{1}{z^2}
  \end{equation}
  As $z \to 0$ in the sparse activation regime, the derivative behaves as $\lim_{z \to 0} | \tilde{\Gamma}'_{FDSA}(z) | = \infty$, proving that the naive FDSA estimator is extremely unstable and has unbounded sensitivity to perturbations in the estimated spectrum.
\end{proof}

\section{Proof of the Spectral Non-Fusion Theorem}
\label{app:proof_non_fusion}
\begin{proof}
  By the Convolution Theorem, pointwise multiplication in the frequency domain is exactly equivalent to circular convolution in the spatial domain:
  \begin{equation}\label{eq:conv_theorem_app}
    \mathcal{F}^{-1}(\Gamma(u, v) \cdot \tilde{X}(u, v)) = G * X
  \end{equation}
  where $G = \mathcal{F}^{-1}(\Gamma) \in \mathbb{R}^{H \times W}$ is the spatial impulse response of the scaling map. Fusing this operation into a preceding convolutional layer with weight $W_{conv}$ of kernel size $K \times K$ (where typically $K=3$ or $K=1$) would require the combined operation to be expressible as a single localized spatial convolution. However, since $\Gamma$ is non-flat, its spatial inverse Fourier transform $G$ has non-zero elements across the entire $H \times W$ grid. The convolution of $W_{conv}$ and $G$ has a spatial support of size $(H + K - 1) \times (W + K - 1)$, which exceeds the local kernel footprint of standard convolution operations. Thus, spectral-domain calibration cannot be folded into standard localized neural network parameters, necessitating active inference hooks.
\end{proof}

\section{Proof of the Phase Coherence and Spatial Edge Preservation Theorem}
\label{app:proof_phase_coherence}
\begin{proof}
  Let $X(y, x)$ and $Y(y, x)$ represent the spatial activation maps of the calibrated merged network and target expert network respectively. By Parseval's energy conservation theorem, the spatial $L_2$ reconstruction error is exactly preserved in the 2D discrete Fourier domain:
  \begin{align}
    D &= \sum_{y=1}^H \sum_{x=1}^W \left| X_{calibrated}(y, x) - Y(y, x) \right|^2 \nonumber \\
    &= \frac{1}{H \cdot W} \sum_{v=1}^H \sum_{u=1}^W \left| \tilde{X}_{calibrated}(u, v) - \tilde{Y}(u, v) \right|^2
  \end{align}
  We substitute the phase-preserving calibrated spectrum $\tilde{X}_{calibrated}(u, v) = \Gamma(u, v) A(u, v) e^{i \phi(u, v)}$ and the target spectrum $\tilde{Y}(u, v) = B(u, v) e^{i \theta(u, v)}$ into the Fourier-domain sum. For ease of notation, let us analyze the term at a single frequency index $(u, v)$:
  \begin{align}
    d(u, v) &= \left| \Gamma(u, v) A(u, v) e^{i \phi(u, v)} - B(u, v) e^{i \theta(u, v)} \right|^2 \nonumber \\
    &= \left( \Gamma(u, v) A(u, v) e^{i \phi(u, v)} - B(u, v) e^{i \theta(u, v)} \right) \left( \Gamma(u, v) A(u, v) e^{-i \phi(u, v)} - B(u, v) e^{-i \theta(u, v)} \right) \nonumber \\
    &= \Gamma(u, v)^2 A(u, v)^2 + B(u, v)^2 - 2 \Gamma(u, v) A(u, v) B(u, v) \cos\left( \phi(u, v) - \theta(u, v) \right)
  \end{align}
  To find the unique minimizer of the reconstruction error $D$ with respect to the phase difference, we must maximize the cosine term $\cos\left( \phi(u, v) - \theta(u, v) \right)$ at all frequencies where $\Gamma(u, v) A(u, v) B(u, v) > 0$. Since the cosine function achieves its unique global maximum of $1$ when its argument is zero, we have:
  \begin{equation}
    \phi(u, v) - \theta(u, v) = 0 \implies \phi(u, v) = \theta(u, v)
  \end{equation}
  This proves that keeping the phase intact is the mathematically unique minimizer of spatial reconstruction error under phase coherence.
  
  Now, suppose there exists a phase perturbation $\Delta \phi(u, v)$ such that the calibrated phase becomes $\phi(u, v) + \Delta \phi(u, v)$. In the spatial domain, a phase shift of $\Delta \phi(u, v) = -u \Delta x - v \Delta y$ corresponds exactly to a spatial translation of the activation map:
  \begin{equation}
    X_{perturbed}(y, x) = X_{calibrated}(y - \Delta y, x - \Delta x)
  \end{equation}
  Under non-linear ReLU activations, spatial translations shift the boundaries of activation regions. Specifically, let $\Omega_1 = \{(y, x) \mid X_{calibrated}(y, x) > 0\}$ and $\Omega_2 = \{(y, x) \mid X_{perturbed}(y, x) > 0\}$ represent the support regions of the unperturbed and perturbed activations. The spatial mismatch between these regions is given by the symmetric difference $\Omega_1 \triangle \Omega_2$. Any non-zero translation $(\Delta x, \Delta y)$ results in a non-empty symmetric difference $\Omega_1 \triangle \Omega_2 \ne \emptyset$, causing spatial feature misalignment, which propagates as localized representation error in subsequent layers. This completes the proof.
\end{proof}

\end{document}
